% Live-book proof adaptation of the arbitrary-ground-field nonreduced
% addendum.  Date: 28 August 2026.

\begin{proof}
The reduction has the same underlying topological space and the same
completed residue fields as \(X\).  Affinoid charts pass to quotients by
nilpotent ideals, so goodness, separatedness, strictness, and pure
dimension one pass to \(X_{\mathrm{red}}\).  The immersion \(\iota\) is
finite and hence has empty relative boundary by
\cite[\S11.4.4]{CLD25}.  The boundary-composition identity
\cite[(11.4.1.3)]{CLD25} therefore shows that
\(X_{\mathrm{red}}\) is boundaryless.  Thus \(X_{\mathrm{red}}\)
satisfies the ambient hypotheses of the reduced theorem
\cref{prop:p21p3-curve-rigidity}.

\smallskip
\emph{Forms, currents, and psh functions.}
By \cref{itm:ch13-pb:red} of
\cref{prop:ch13-pullback-properties}, proved from
\cite[Proposition~15.1.3(b) and the following form variant]{CLD25},
pullback gives an isomorphism

\begin{equation}
\label{eq:p21p3-nr-form-isomorphism}
       \iota^*:A_X^{p,q}\xrightarrow{\ \sim\ }
                       A_{X_{\mathrm{red}}}^{p,q}.
\end{equation}

The point needed for currents is topological, not merely algebraic.  A
datum for the compact-test topology consists of a compact support and a
finite family of moment charts.  Restriction sends such a datum on \(X\)
to the same datum on \(X_{\mathrm{red}}\).  Conversely, the cited
almost-good result applies because \(X\) is good.  At every point of the
compact support, a finite tuple of invertible functions defining a reduced
moment chart lifts through the nilpotent quotient to a tuple of unit germs
on \(X\).  Each lift has the same absolute value as its reduction at every
analytic point.  After shrinking the ordinary open neighbourhoods and
using compactness, one obtains a finite lifted refinement; product moment
maps give a common refinement on overlaps.

Restriction and these finite lifted refinements are cofinal systems of
compact chart data for the final locally convex topologies.  On a lifted
chart, the tropical images, Lagerberg representatives, polyhedral
refinements, and every \(C^r\)-seminorm are identical.  Consequently

\begin{equation}
\label{eq:p21p3-nr-test-homeomorphism}
 \iota^*:A^{p,q}_{X,c}\xrightarrow{\ \sim\ }
                         A^{p,q}_{X_{\mathrm{red}},c}
\end{equation}

is a homeomorphism.  The same chart description proves compatibility with
wedge product, \(J\), the differentials, supports, and all positivity cones
of forms.  Taking continuous duals gives

\begin{equation}
\label{eq:p21p3-nr-current-isomorphism}
 \iota_*:D^{a,b}(X_{\mathrm{red}})
              \xrightarrow{\ \sim\ }D^{a,b}(X),
 \qquad
 \langle\iota_*S,\alpha\rangle
       =\langle S,\iota^*\alpha\rangle .
\end{equation}

This isomorphism identifies the weak topologies.  It preserves positivity
because the test-form cones agree, and it preserves closedness because
\(\iota^*\) commutes with \(\mathrm d'\) and \(\mathrm d''\).
Multiplication by a smooth function commutes with \(\iota_*\) directly by
transposition.  If \(S\) is positive, the Borel product of
\cite[Example~7.2.12]{CLD25} is defined through the measures
\(\mu_{S,\omega}\) attached to positive test forms \(\omega\).  Under
\(\iota^*\), the positive test forms and these attached measures agree.
Thus \(\iota_*(fS)=f\iota_*S\) for every locally bounded Borel function
\(f\).

If \(S=S_1-S_2\) locally with \(S_1,S_2\) positive, define
\(fS=fS_1-fS_2\).  This is independent of the decomposition.  Indeed, from
\(S_1-S_2=S'_1-S'_2\) one gets
\(S_1+S'_2=S'_1+S_2\); equality of these positive currents gives equality
of all their attached measures, and hence
\(fS_1+fS'_2=fS'_1+fS_2\).  The local definitions glue, and the same
calculation proves compatibility with \(\iota_*\).  Nothing in this
argument defines \(fS\) for an arbitrary order-zero current.

For the psh assertion, stage zero of the CLD filtration is the same on the
two spaces: smoothness and positivity of
\(\mathrm d'\mathrm d''u\) agree by
\eqref{eq:p21p3-nr-form-isomorphism}.  Successor stages take decreasing
limits, the continuous filtration takes local uniform limits, and limit
ordinal stages take local unions.  The underlying open sets and functions
are identical.  Transfinite induction in
\cref{def:ch16-psh-filtration} therefore identifies the psh and continuous
psh sheaves at every stage and after stabilization.  The \(G\)-form
argument is the same, using the \(G\)-form variant immediately after
\cite[Proposition~15.1.3]{CLD25}.  This proves
\cref{itm:p21p3-nr:calculus}.

\smallskip
\emph{The integration current and generic multiplicity.}
By \cite[Definition~13.2.5 and \S13.5.2]{CLD25}, the fundamental cycles
are

\[
                 [X]=\sum_i m_i[X_i],
                 \qquad
                 [X_{\mathrm{red}}]=\sum_i[X_i].
\]

Let
\(\bar\nu_i:\widetilde X_i\to X_i\hookrightarrow X_{\mathrm{red}}\),
so that \(\nu_i=\iota\circ\bar\nu_i\).  The normalization is finite and
generically an isomorphism.  The compatibility of integration with
analytic cycles in \cref{thm:current-fundamental-cycle}, namely
\cite[Proposition~16.6.3]{CLD25}, identifies the current of \([X_i]\) with
\((\bar\nu_i)_*\delta_{\widetilde X_i}\) on \(X_{\mathrm{red}}\), and
with \((\nu_i)_*\delta_{\widetilde X_i}\) after pushforward to \(X\).
It follows that

\[
 \delta_X=\sum_i m_i\delta_i,
 \qquad
 \iota_*\delta_{X_{\mathrm{red}}}=\sum_i\delta_i.
\]

The irreducible components are locally finite by
\cite[\S11.2.5]{CLD25}, so both sums are locally finite.  Their equality
forces \(m_i=1\) for every \(i\): restrict to a nonempty normal open subset
of \(X_i\) meeting no other component and use the nonzero integration
current there.  Conversely, \(m_i=1\) for every \(i\) visibly gives
equality.  Since the length of the generic local ring along \(X_i\) is
\(m_i\), this is equivalent to generic reducedness along every component.
This proves \cref{itm:p21p3-nr:integration}.

\smallskip
\emph{Classification.}
Let \(T\) be positive and closed on \(X\).  By
\eqref{eq:p21p3-nr-current-isomorphism}, there is a unique positive
closed current \(T_{\mathrm{red}}\) on \(X_{\mathrm{red}}\) with
\(T=\iota_*T_{\mathrm{red}}\).  The reduced classification
\cref{itm:p21p3-curve:classification} of
\cref{prop:p21p3-curve-rigidity} gives

\[
                  T_{\mathrm{red}}=\sum_i c_i
                    (\bar\nu_i)_*\delta_{\widetilde X_i},
                  \qquad c_i\geq0.
\]

Pushing forward by \(\iota\) proves
\eqref{eq:p21p3-nr-classification}, including local finiteness and
uniqueness.  If \(X_{\mathrm{red}}\) is connected and rig-smooth, its
regular local rings are domains, so its irreducible components are
disjoint open-and-closed subsets.  It therefore has one component, with
one generic multiplicity \(m\).  Then
\(\delta_X=m\delta_1\), and
\(c_1\delta_1=(c_1/m)\delta_X\).  This proves
\cref{itm:p21p3-nr:classification}.

\smallskip
\emph{The strict-rigid exceptional locus.}
Let
\(\bar\nu:\coprod_i\widetilde X_i\to X_{\mathrm{red}}\) be the
normalization map.  On a good affinoid domain put

\[
 \mathcal Q:=\bar\nu_*\mathcal O_{\coprod_i\widetilde X_i}
                 /\mathcal O_{X_{\mathrm{red}}},
 \qquad
 \mathfrak c:=\operatorname{Ann}_{\mathcal O_{X_{\mathrm{red}}}}
                         (\mathcal Q).
\]

Let \(Z\) be the union of the conductor locus

\[
 \operatorname{Supp}(\mathcal Q)
   =V(\mathfrak c)
   =\operatorname{Supp}(\mathcal O_{X_{\mathrm{red}}}/\mathfrak c)
\]

and the pairwise intersections of distinct irreducible components.
Analytic local rings on a good space are excellent by
\cite[\S11.1.10]{CLD25}; normalization is finite and its components
correspond to the irreducible components by
\cite[\S11.3.1]{CLD25}.  Hence \(\mathcal Q\) and the conductor are
coherent.  Their support misses every generic point of the reduced pure
curve and is zero-dimensional.  The same dimension argument applies to
the proper intersections of distinct components.  Local finiteness of the
components, again \cite[\S11.2.5]{CLD25}, shows that \(Z\) is closed and
locally finite.

Here strictness replaces algebraic closedness.  Locally, a
zero-dimensional closed analytic subspace of a strict affinoid domain is
the spectrum of a zero-dimensional strict affinoid algebra \(B\).  By
\cite[Theorem~4.5]{Tate71Rigid}, the algebra \(B\) is noetherian; having
Krull dimension zero, it is therefore Artinian.
For a point \(z\) of its Berkovich spectrum, the kernel of the associated
seminorm is maximal; the strict affinoid Nullstellensatz
\cite[Theorem~4.5]{Tate71Rigid} makes
\(B/\ker(z)\) finite over \(K\).  This finite-dimensional \(K\)-algebra is
already complete, so
\(\mathcal H(z)=B/\ker(z)\) is finite over \(K\).  Thus every point of
\(Z\) is rigid, although it need not be \(K\)-rational.

Put

\begin{equation}
\label{eq:p21p3-nr-conductor-preimage}
                         Z_i:=\bar\nu_i^{-1}(Z).
\end{equation}

A finite morphism induces finite extensions of completed residue fields,
so every point of \(Z_i\) is rigid.  A finite algebra over a strict
affinoid algebra is strict by \cite[Theorem~4.4]{Tate71Rigid};
consequently each \(\widetilde X_i\) is strict.  The finite maps
\(X_i\to X_{\mathrm{red}}\) and
\(\widetilde X_i\to X_i\) have empty relative boundary by
\cite[\S11.4.4]{CLD25}.  Applying
\cite[(11.4.1.3)]{CLD25} twice shows that every
\(\widetilde X_i\) is boundaryless.  It is good and separated, and it is
rig-smooth by hypothesis.  Away from \(Z\), the component images are
disjoint and each \(\bar\nu_i\) is an analytic isomorphism.  These are
exactly the hypotheses needed below for the rigid zero-atom input; no
assertion about arbitrary type-\(1\) points is used.

\smallskip
\emph{The multiplicity-weighted curvature identity.}
Fix \(W,U,V\) as in \cref{itm:p21p3-nr:bedford-taylor}, identify \(W\)
with the corresponding open subset of \(X_{\mathrm{red}}\), and put
\(\widetilde W_i=\bar\nu_i^{-1}(W)\).  Denote the base-changed maps by

\[
 \nu_{i,W}^{X}:\widetilde W_i\longrightarrow W\subset X,
 \qquad
 \nu_{i,W}^{\mathrm{red}}:\widetilde W_i\longrightarrow
                         W\subset X_{\mathrm{red}}.
\]

Pullback preserves continuous psh functions by
\cite[Proposition~8.1.5]{CLD25}.  On \(\widetilde W_i\), define positive
Radon measures

\[
 \mu_i^U:=\mathrm d'\mathrm d''[\bar\nu_i^*U],
 \qquad
 \mu_i^V:=\mathrm d'\mathrm d''[\bar\nu_i^*V],
 \qquad
 \sigma_i:=\mu_i^U-\mu_i^V.
\]

Their positivity is \cite[Lemma~8.5.5]{CLD25}.  For a continuous function
\(F\) on \(W\), multiply the two locally finite identities in
\eqref{eq:p21p3-nr-fundamental-cycle} by \(F\), using
\cite[Example~7.2.12]{CLD25}.  Evaluation on a compactly supported test
form and the defining transpose pairing for finite pushforward give

\begin{align}
 [F]_X
   &=\sum_i m_i(\nu_{i,W}^{X})_*
       \bigl(((\nu_{i,W}^{X})^*F)\delta_{\widetilde W_i}\bigr),
 \label{eq:p21p3-nr-function-current-X}\\
 [F]_{X_{\mathrm{red}}}
   &=\sum_i(\nu_{i,W}^{\mathrm{red}})_*
       \bigl(((\nu_{i,W}^{\mathrm{red}})^*F)
                    \delta_{\widetilde W_i}\bigr).
 \label{eq:p21p3-nr-function-current-red}
\end{align}

For \(F=U,V\), differentiate these identities and use the continuous-psh
projection formula \cite[Proposition~8.5.6]{CLD25}; equivalently, proper
pushforward commutes with the current differentials by transposition.
Thus equality of the curvature currents computed with \(\delta_X\) is
equivalent to

\begin{equation}
\label{eq:p21p3-nr-weighted-measure-sum}
                       \sum_i m_i(\nu_{i,W}^{X})_*\sigma_i=0,
\end{equation}

whereas equality computed with \(\delta_{X_{\mathrm{red}}}\) is
equivalent to

\begin{equation}
\label{eq:p21p3-nr-unweighted-measure-sum}
                       \sum_i
                         (\nu_{i,W}^{\mathrm{red}})_*\sigma_i=0.
\end{equation}

Write \(Z_W=Z\cap W\) and
\(Z_{i,W}=Z_i\cap\widetilde W_i=\bar\nu_i^{-1}(Z_W)\).  Away from \(Z_W\),
the component images are disjoint and the relevant normalization map is
an isomorphism.  Restricting either
\eqref{eq:p21p3-nr-weighted-measure-sum} or
\eqref{eq:p21p3-nr-unweighted-measure-sum} to such an open set, and
using \(m_i>0\) in the weighted case, yields

\begin{equation}
\label{eq:p21p3-nr-component-support}
                       \operatorname{Supp}(\sigma_i)
                       \subseteq Z_{i,W}
                       \quad\text{for every }i.
\end{equation}

Every point of \(Z_{i,W}\) is rigid.  Around such a point, choose an open
neighbourhood contained in \(\widetilde W_i\).  It is a good separated
strict rig-smooth boundaryless curve, so
\cref{lem:p21p3-gfc-rigid-zero-atom} applies to the restrictions of
\(\bar\nu_i^*U\) and \(\bar\nu_i^*V\).  Neither positive measure
\(\mu_i^U\) nor \(\mu_i^V\) has an atom there.  A signed Radon measure
supported on a locally finite set is the locally finite sum of its point
masses.  Hence \eqref{eq:p21p3-nr-component-support} implies

\begin{equation}
\label{eq:p21p3-nr-component-measure-equality}
                         \sigma_i=0\qquad\text{for every }i.
\end{equation}

This argument works starting from either the weighted or the unweighted
identity.  Conversely,
\eqref{eq:p21p3-nr-component-measure-equality} implies both identities.
This proves the exact curvature equivalence
\eqref{eq:p21p3-nr-curvature-equivalence}; in particular, generic
multiplicities cannot conceal cancellation at a component intersection.

Finally put \(h=U-V\) and write \(T=\sum_i c_i\delta_i\).  Multiplication
by \(h\) is defined for the positive current \(T\) by
\cite[Example~7.2.12]{CLD25}; continuity makes \(h\) locally bounded.
Proper pushforward, linearity of the current differentials, and
\eqref{eq:p21p3-nr-component-measure-equality} give

\begin{align*}
 \mathrm d'\mathrm d''(hT|_W)
 &=\sum_i c_i(\nu_{i,W}^{X})_*
     \mathrm d'\mathrm d''
       [\bar\nu_i^*U-\bar\nu_i^*V]\\
 &=\sum_i c_i(\nu_{i,W}^{X})_*(\mu_i^U-\mu_i^V)=0.
\end{align*}

This is \eqref{eq:p21p3-nr-bt-conclusion}.  It proves precisely the
Bedford--Taylor independence identity for the continuous psh pairs stated
in the proposition.  The proof does not invoke, and the proposition does
not assert, the additional discontinuous-psh quantifier in
\cite[Definition~8.6.2]{CLD25}.  This completes the proof.
\end{proof}
